\section{Compactness, vocabulary and spaces of theories}
\label{sec:p03-compactness-vocabulary}

Compactness bounds the information needed to distinguish structures.
It also has a topological formulation, but the fixed-cardinal versions
must distinguish single sentences from whole theories placed in blocks.

\begin{hbdefinition}{def:p03-cv-cardinals}{Two cardinal parameters}
Let \(\kappa\geq\lambda\geq\aleph_0\). A logic is
\emph{\((\kappa,\lambda)\)-compact} if, for every vocabulary and
every theory \(T\) with \(|T|\leq\kappa\), satisfiability of
every subset of \(T\) of size less than \(\lambda\) implies
satisfiability of \(T\). Thus \((\kappa,\aleph_0)\)-compactness
is the \(\kappa\)-compactness of
\cref{def:p03-w03-p03-041-compactness}.
The cardinality bound on the whole theory is non-strict; the bound
on its tested subtheories is strict
\cite[Ch.~II, Definition~5.1.1, p.~59]{Ebbinghaus1985ExtendedLogics}.
\end{hbdefinition}

\subsection{The essential vocabulary of a sentence}

\begin{hbdefinition}{def:p03-cv-support}{Semantic support}
A sentence \(\varphi\in\mathscr L[\tau]\) \emph{depends only on}
a sub-vocabulary \(\sigma\subseteq\tau\) if any two
\(\tau\)-structures whose \(\sigma\)-reducts are isomorphic agree
on \(\varphi\). Such a \(\sigma\) is a semantic support.
This says something about satisfaction; it does not assert that the
same written sentence belongs to \(\mathscr L[\sigma]\).
Sort symbols and the sorts attached to retained symbols are included
when counting a vocabulary.
\end{hbdefinition}

\begin{hbtheorem}{thm:p03-cv-support}{proved=annex:p03.cv.1}{Compactness bounds semantic support}
Let \(\mathscr L\) satisfy the general framework of
\cref{def:p03-w03-p03-047-ebbinghaus-logic} and the atom, Boolean
and particularization properties of
\cref{def:p03-w03-p03-047-closure}. It may be many-sorted.
If it is \((\kappa,\lambda)\)-compact and \(|\tau|\leq\kappa\),
every sentence in \(\mathscr L[\tau]\) has a semantic support
of cardinality less than \(\lambda\).
In particular every sentence of a compact logic with these properties
has finite semantic support. This is the support conclusion of
\cite[Ch.~II, Proposition~5.1.2, p.~60]{Ebbinghaus1985ExtendedLogics},
with the sort identifications made explicit.
\end{hbtheorem}

\begin{hbcorollary}{thm:p03-cv-support-translation}{proved=annex:p03.cv.2}{A sentence over its semantic support}
Let \(\mathscr L\) be one-sorted and regular. If
\(\varphi\in\mathscr L[\tau]\) has semantic support
\(\sigma\subseteq\tau\), then some
\(\psi\in\mathscr L[\sigma]\) is equivalent to \(\varphi\)
over all \(\tau\)-structures. Consequently a compact one-sorted
regular logic has an equivalent finite-vocabulary sentence for every
sentence, even when its syntax does not require finite occurrence.
\end{hbcorollary}

\begin{hbexample}{def:p03-cv-unbounded-support}{No uniform support bound for the infinitary language}
For an arbitrary infinite set \(J\), take unary predicates \(P_j\)
and the \(L_{\infty\omega}\)-sentence
\[
 \bigwedge_{j\in J}\forall x\,P_j(x).
\]
Every semantic support must contain every \(P_j\): on a one-element
domain, changing just an omitted predicate from full to empty changes
the sentence's truth without changing the supported reduct.
As \(|J|\) is arbitrary, no single cardinal bounds the essential
vocabularies of all sentences of this language.
\end{hbexample}

\subsection{Spaces of realized complete theories}

\begin{hbdefinition}{def:p03-cv-space}{The theory space of a vocabulary}
Fix \(\tau\) such that \(\mathscr L[\tau]\) is a set, and assume
Boolean closure and first-order tautologies. Define
\[
 \begin{aligned}
 X_{\mathscr L,\tau}
   &=\{\operatorname{Th}_{\mathscr L}(\mathfrak A):
                    \mathfrak A\in\operatorname{Str}[\tau]\},\\
 U_\varphi&=\{p\in X_{\mathscr L,\tau}:\varphi\in p\}.
 \end{aligned}
\]
This is a set: it is a subclass of the power set of
\(\mathscr L[\tau]\). Give it the topology generated by the
\(U_\varphi\). Its points are complete theories actually satisfied
by structures, not arbitrary unfulfilled truth assignments.
Equivalently, choose one structure from each logical-equivalence
class and use the sentence-defined model classes as basic opens.
The two presentations identify corresponding points and basic opens
\cite[Ch.~II, \S5.1, p.~59]{Ebbinghaus1985ExtendedLogics}.
\end{hbdefinition}

\begin{hbtheorem}{thm:p03-cv-space}{proved=annex:p03.cv.3}{Logical and topological compactness}
The sets \(U_\varphi\) form a clopen basis and
\(X_{\mathscr L,\tau}\) is Hausdorff. It is compact if and only if
every finitely satisfiable theory in vocabulary \(\tau\) is
satisfiable. Consequently, for a small logic (each
\(\mathscr L[\tau]\) a set), full logical compactness is equivalent
to compactness of every such theory space.
This extends the first-order Stone-space argument used in
\cref{thm:p03-robinson-joint-consistency}; first-order type-space
compactness itself is \cref{thm:type-space-compact}.
\end{hbtheorem}

\begin{hbdefinition}{def:p03-cv-blocks}{Topological and block compactness}
A space is \emph{\((\kappa,\lambda)\)-compact} if every family
of at most \(\kappa\) closed sets with empty intersection has a
subfamily of size less than \(\lambda\) with empty intersection.
For the fixed vocabulary \(\tau\), the corresponding
\emph{block condition} says: for every family \((T_i)_{i\in I}\)
of theories with \(|I|\leq\kappa\), if
\(\bigcup_{i\in J}T_i\) is satisfiable for each \(J\subseteq I\)
of size less than \(\lambda\), then
\(\bigcup_{i\in I}T_i\) is satisfiable. Each block \(T_i\) is an
arbitrary subset of \(\mathscr L[\tau]\); there is no bound of
\(\kappa\) on the size of a block or of their union.
\end{hbdefinition}

\begin{hbtheorem}{thm:p03-cv-blocks}{proved=annex:p03.cv.4}{Closed sets and theory blocks}
The theory space \(X_{\mathscr L,\tau}\) is topologically
\((\kappa,\lambda)\)-compact exactly when the block condition
holds. Ordinary logical \((\kappa,\lambda)\)-compactness tests
single sentences, equivalently families of sentence-defined basic
clopens; it need not imply this condition for arbitrary closed sets.
If \(|\mathscr L[\tau]|\leq\kappa\), however, logical
\((\kappa,\lambda)\)-compactness does imply the block condition
at this vocabulary.
\end{hbtheorem}

\begin{hbtheorem}{thm:p03-cv-qone-space}{proved=annex:p03.cv.5}{A compactness boundary}
The logic \(\operatorname{FO}(Q_1)\) is countably compact, but
for the vocabulary \(\{c_\alpha:\alpha<\omega_1\}\) its space
of realized complete theories is not topologically
\((\aleph_0,\aleph_0)\)-compact. The latter failure is witnessed
by countably many closed sets whose every finite subfamily even
has a point represented by a finite structure.
No continuum hypothesis is assumed.
\end{hbtheorem}

Thus the unrestricted correspondence in \cref{thm:p03-cv-space}
does not transfer unchanged to fixed cardinal bounds. The extra
strength lies in treating an entire theory as one closed-set block,
as in the strengthened compactness discussed in
\cite[Ch.~II, \S5.1, p.~60]{Ebbinghaus1985ExtendedLogics}.
